% =============================================================================
% CAGE — Introduction: pattern-matched options (PUC ABNT formatting)
% Copy-paste the blocks you want into my-cap1.tex.
%
% RULES FOLLOWED (per /Users/lucasmariano/CAGE-Formats):
%   - \cite{key} at the END of a sentence; \citeonline{key} when the author is in
%     the prose; enumerations via compactitem with \item[a)]; tables carry the
%     mandatory \captionfont{...Fonte...} line.
%   - Every \cite key below is ALREADY in article-references.bib (verified).
%   - Each research question / hypothesis states its GROUNDING IN CAGE and the
%     PHASE that will answer it (Phase 1 done; Phases 2-3 future).
%
% PROVENANCE (anti-hallucination): [data] = Phase-1 results; [cite] = verified
% reference; numbers (37.4%, 65.7%, 70.4%, 11.6%, 7.6x, 0.570->0.504, 40 GB) are
% all verified. Nothing about Phase 2/3 is asserted as a result.
% =============================================================================


% =============================================================================
% 1.3  QUESTÕES DE PESQUISA E HIPÓTESES  (the centerpiece — paste after the
%      central question already in my-cap1.tex, line ~18)
% =============================================================================

The central question decomposes into five research questions, each anchored in a
component of CAGE and each scheduled to be answered as the framework advances
through its phases.

\begin{compactitem}
  \item[\textbf{RQ1}] What metric suite jointly captures serving efficiency, cache-reuse behaviour, and answer faithfulness, so that the trade-off is measured rather than assumed?
  \item[\textbf{RQ2}] Under controlled conditions, what is the measurable effect of cache reuse on serving performance and on answer faithfulness, relative to retrieval-backed and hybrid baselines?
  \item[\textbf{RQ3}] As serving scales to GPU and multi-node high-performance-computing (HPC) deployments under heavy KV-cache memory pressure, how does distributed KV-cache efficiency trade off against information-retrieval quality?
  \item[\textbf{RQ4}] Which telemetry signals are necessary and sufficient to attribute serving behaviour to its true cause --- local reuse, retrieval success, or cross-node transfer?
  \item[\textbf{RQ5}] Does compressing each paradigm's native object --- the retrieved text for RAG and the resident KV cache for CAG --- reduce serving cost without re-ordering the answer-quality ranking between caching and retrieval, and at what compression level does aggressive reduction begin to cost faithfulness?
\end{compactitem}

These questions are tested through six hypotheses. Each is falsifiable, grounded in
a specific component of the framework and tied to the phase in which CAGE will
test it; together they translate the central question into concrete, measurable
claims about the efficiency--quality trade-off.

\begin{compactitem}
  \item[\textbf{H1}] Local prompt-cache reuse --- not retrieval success --- is the primary driver of Time-To-First-Token (TTFT) reduction. This claim is foundational to the central question: it asserts that the efficiency of cache-aware serving comes from \textit{reuse}, which is precisely why the framework must instrument reuse separately from retrieval success. \textit{Grounding:} CAGE's cache-hit and prompt-cached telemetry. \textit{Evaluated in:} Phases 1--2, under real prefill cost.
  \item[\textbf{H2}] When the gold context is already available, retrieval lowers answer faithfulness, and this loss is visible to a strict Natural-Language-Inference (NLI) metric but invisible to soft-matching similarity. It addresses the quality side of the trade-off and justifies the framework's reliance on strict faithfulness: an efficiency choice must not hide a grounding cost that the metric cannot detect. \textit{Grounding:} CAGE's NLI faithfulness metric, with BERTScore retained as a negative control~\cite{espejel2023ragas, ares2024}. \textit{Evaluated in:} Phases 1--2.
  \item[\textbf{H3}] In distributed deployments, the efficiency gained by reusing and distributing the KV cache is bounded by cross-node transfer and cold-start costs, producing a measurable efficiency--quality frontier. This is the central trade-off named in the title --- distributed KV-cache efficiency against retrieval quality --- and locating that frontier is the purpose of the HPC study. \textit{Grounding:} CAGE's distributed baseline and its cross-node transfer telemetry~\cite{zhong2024distserve, qin2024mooncake}. \textit{Evaluated in:} Phase 3.
  \item[\textbf{H4}] Near-lossless KV-cache compression (FP8 quantisation, or low-rank attention such as MLA) holds faithfulness at parity while postponing the eviction/saturation point, whereas aggressive eviction under heavy memory pressure degrades it. Because memory pressure is what forces compression at scale, this hypothesis links the dominant efficiency lever directly to answer quality --- the exact question CAGE poses under HPC load. \textit{Grounding:} the \texttt{compressed\_cag} arm of CAGE's compression axis and its KV-bytes-per-token telemetry~\cite{deepseekv2, kvtuner, chen2025pitfalls}. \textit{Evaluated in:} Phases 2--3.
  \item[\textbf{H5}] Prompt compression lowers the per-request cost of retrieval but does not close RAG's grounding gap; compression is therefore largely orthogonal to the CAG-versus-RAG quality ordering. It tests whether cheaper retrieval can let RAG catch up to caching on quality --- a direct probe of the paradigm comparison at the heart of the work. \textit{Grounding:} the \texttt{compressed\_rag} arm of the compression axis, using LLMLingua-2 prompt compression~\cite{llmlingua2}. \textit{Evaluated in:} Phase 2.
  \item[\textbf{H6}] Tail latency (p95), not the mean, governs the user experience of cache-aware serving, and cold replicas dominate the tail. This justifies the framework's percentile reporting and is decisive in the distributed regime, where cold caches create the tails that define perceived efficiency. \textit{Grounding:} CAGE's percentile (p50/p95) latency reporting. \textit{Evaluated in:} Phases 1--3.
\end{compactitem}

The compression axis warrants explicit attention, because it is where the
efficiency--quality trade-off becomes unavoidable at scale. Production serving rarely
runs at full precision or full context: prompts are compressed and KV caches are
quantised or evicted to fit memory and budget~\cite{kvcompresssurvey, li2024kvsurvey}.
Yet these techniques are evaluated almost exclusively for throughput and memory; their
effect on answer faithfulness is under-studied, and recent work shows that KV-cache
compression can silently cause a model to ignore instructions --- a ``selective
amnesia'' --- on realistic prompts~\cite{chen2025pitfalls}. Crucially, the two
paradigms compress \textit{different objects}: prompt tokens for RAG (e.g.\
LLMLingua-2~\cite{llmlingua2}) versus cached tensors for CAG (e.g.\ FP8 or Multi-head
Latent Attention~\cite{deepseekv2}). No prior study crosses compression with the
cache-versus-retrieve choice while measuring grounding; this is the gap that RQ5
validates. CAGE closes it with an explicit 2$\times$2 axis (\texttt{compressed\_rag}
$\times$ \texttt{compressed\_cag}) that pairs a token-retention ratio and a
KV-bytes-per-token measure with the same faithfulness metrics, so that every
efficiency gain is reported against its quality cost; H4 and H5 state the expected
outcomes.

Table~\ref{tab:rq_grounding} summarises this mapping. \textit{Why this table:} it makes
the research design traceable and falsifiable at a glance --- every question and
hypothesis is anchored to a concrete CAGE component and scheduled to a phase, so a
reader can see that no claim is left unjustified or unscheduled.

\begin{table}[H]
\centering
\setlength{\abovecaptionskip}{0pt}
\caption[Research questions and hypotheses, their grounding, and evaluating phase]{Mapping of research questions and hypotheses to their grounding in CAGE and the phase in which each is evaluated.}
\label{tab:rq_grounding}
\footnotesize
\begin{tabular}{@{}p{0.16\textwidth}p{0.50\textwidth}p{0.24\textwidth}@{}}
\toprule
\textbf{Item} & \textbf{Grounding in CAGE (justification)} & \textbf{Evaluated in} \\
\midrule
RQ1            & Layered telemetry + semantic-quality metric suite & Phase 1 \\ \addlinespace
RQ2, RQ3, H1  & Seven-baseline taxonomy + prompt-cache telemetry  & Phases 1--2 \\ \addlinespace
RQ4 (quality), H2 & NLI faithfulness + LettuceDetect grounding + BERTScore control & Phase 1 \\ \addlinespace
H6            & Tail-latency (p50/p95) reporting & Phase 1 \\ \addlinespace
RQ5, H4, H5   & The 2$\times$2 compression axis (\texttt{compressed\_cag} / \texttt{compressed\_rag}), with token-retention and KV-bytes/token metrics & Phase 2 \\ \addlinespace
RQ3, RQ4 (transfer), H3 & Distributed baseline + cross-node KV-transfer telemetry & Phase 3 \\
\bottomrule
\end{tabular}
\captionfont{\small{\textbf{\\Source: Prepared by the author.}}}
\end{table}


% =============================================================================
% 1.1  MOTIVAÇÃO  (alternative / additional opening — OPTION A is quantitative)
% =============================================================================

% --- OPTION A (KV-cache-as-bottleneck; quantitative) ---
Large Language Models have made inference, rather than training, the dominant
operational cost of deployed AI, and within inference the Key-Value (KV) cache has
become the binding constraint. The cache grows linearly with context length and
batch size, and at long context it can equal the model weights themselves: for
Llama-3.1-70B at a 128K-token context, a single sequence's KV cache reaches roughly
40~GB~\cite{li2024kvsurvey}. Two designs respond to this pressure in opposite ways:
Retrieval-Augmented Generation (RAG) fetches fresh context per request, while
Cache-Augmented Generation (CAG) reuses previously processed context and its cached
state~\cite{yu2024dontdorag}. Which bet is better, and at what cost to answer
quality, is precisely what current practice cannot say.

% --- OPTION B (evaluation-gap; shorter) ---
% The shift from single-turn retrieval toward reuse-heavy serving has outpaced the
% tools used to evaluate it: the community measures retrieval pipelines with quality
% frameworks and serving systems with throughput benchmarks, but no methodology
% measures both at once --- even though the value of cache reuse is exactly a joint
% question of efficiency and grounding.


% =============================================================================
% 1.5  JUSTIFICATIVA
% =============================================================================

% --- OPTION A (the gap) ---
The justification is a measurable gap. Quality-oriented evaluation frameworks
instrument the retriever and generator but not the serving
layer~\cite{espejel2023ragas, ares2024, li2024ragbench}, whereas systems-oriented
serving research optimises goodput and cache disaggregation but does not score
answer grounding~\cite{zhong2024distserve, qin2024mooncake}. To our knowledge, no
framework jointly instruments cache reuse, serving efficiency, and grounding under a
common protocol --- which is exactly the evidence a practitioner needs to choose
between caching and retrieval.

% --- OPTION B (timeliness + payoff) ---
% The work is timely and practical: KV-cache reduction --- quantisation, eviction,
% low-rank attention --- is one of the most active frontiers of LLM
% serving~\cite{li2024kvsurvey}, because at long context the cache can equal the
% model weights. Yet these advances are evaluated for speed and memory, almost never
% for their effect on answer faithfulness. A framework that ties efficiency to
% grounding gives architects a defensible basis for deciding when cache reuse,
% retrieval, or a hybrid policy is preferable at a target quality.


% =============================================================================
% 1.6  CONTRIBUIÇÕES
% =============================================================================

The main contributions are:

\begin{compactitem}
  \item[a)] CAGE, a layered and modular framework that jointly evaluates serving performance, cache telemetry, and semantic quality across cache-backed, retrieval-backed, hybrid, and distributed baselines under common conditions;
  \item[b)] a controlled local validation that establishes the framework's metric and orchestration integrity and yields a first quantitative characterisation of the reuse-versus-retrieval trade-off, including a soft-matching metric shown to be a non-discriminative negative control;
  \item[c)] a research plan and methodology for quantifying distributed KV-cache efficiency against information-retrieval quality under HPC workloads; and
  \item[d)] a reproducible software and serving-telemetry stack that captures signals --- speculative-decode acceptance, KV-compression ratio, and prompt-token source --- that request-level metrics do not expose.
\end{compactitem}


% =============================================================================
% 1.7  LIMITAÇÕES DA PESQUISA  (recommended; matches the PUC pattern)
% =============================================================================

This work is bounded in four ways, stated explicitly.

\begin{compactitem}
  \item[a)] The completed validation runs on a CPU backend that does not induce the extreme KV-cache memory pressure that motivates CAG; its absolute latency values are not generalisable, and only relative orderings should be read from it.
  \item[b)] The Phase-1 distributed baseline performs orchestration-level routing with a modelled transfer cost; it does not move real KV tensors between nodes --- real cross-node transfer is a Phase-3 target.
  \item[c)] The Phase-1 quality constants are preliminary; the faithfulness and similarity metrics were subsequently strengthened, and the numbers are re-established under the corrected protocol in Phase 2.
  \item[d)] The local validation uses a single dataset (SQuADv2) and model family (Qwen3-4B); dataset and model diversity are introduced in the GPU phases.
\end{compactitem}


% =============================================================================
% NOTE — table compliance fix for my-cap1.tex
% Both tables in my-cap1.tex (tab:rag_vs_cag and tab:related_work) are missing the
% mandatory ABNT "Fonte:" line. Add the line below just before each table closes
% (now applied to both my-cap1.tex tables):
%   \captionfont{\small{\textbf{\\Source: Prepared by the author.}}}
% =============================================================================
